% H50: field vs coupling from transfer lags. Sources: hypotheses/H50-field-vs-coupling-transfer-lag/{README.md, summary/content.tex, summary/meta.json};
% writeup/visuals/H50-field-vs-coupling-transfer-lag/{caption.tex, README.md}; interpretation/swarm-constants.json (J1 named vs unnamed).
\section{H50: Activity moves with the scheduler; talk moves one read-out call at a time}

\textbf{The question:}
When agents co-move, is a common input hitting all of them (a field), or do they read each other (a coupling)?
Can the lag between input and response tell the two apart?

\textbf{The intuition:}
Agents act only at model calls.
A field, for example the runner's daily start, reaches every agent at its next call.
So all agents respond within about one call cycle: the field has zero relative lag.
A coupling must pass through a message.
The message enters a recipient's context only at its next call, the read-out call.
So coupling lags come in whole call cycles, or hops.
Hop 0 is the call already in flight, which cannot see the message; hop 1 is the read-out call.
The test is a regression discontinuity: it compares calls that start just after a message with calls that start just before it.
A field gives no step there; a coupling gives a step at hop 1.

\textbf{What we measured:}
The read-out jump $J_1$ is the extra chance that a recipient's hop-1 call is a talk call, minus a shifted-time placebo.
It uses the DQ1 context ledger, which records which call received which message.
The analysis also fits impulse responses (the response over time to one input) for day starts, pauses, human messages, nudges and platform errors.
It then splits co-movement into a field part and a coupling part.
It covers 71 units in 35 goal periods outside the reserved data, in all three regimes, plus 4 period-native tests.
\emph{Scheduler field:} measured as an input against a shifted-input null, with an all-running-window variant.
\emph{External drive:} human messages, nudges, operator bookends and errors enter as fields.
Goal kickoffs are not among the inputs, so this impostor is only partly removed.
\emph{Shared model prior:} not relevant, because a family prior cannot make a step at the read-out call.
\emph{Common input:} removed by the comparison of the in-flight call with the read-out call; the in-flight side is not raised.
\emph{Synthetic check:} 15 planted scenarios $\times$ 5 seeds on real village schedules.
$J_1$ recovers planted coupling with a bias of $-2$ to $-17\%$ in regime III.
It reads zero for fields and slow common drives, and it finds a planted 2-hop dead time as onset at hop 3.
It goes negative for a coupling that acts within the call in flight.
False positives occur in 3/35 null seeds.

\begin{figure}[h]
\includegraphics[width=\textwidth]{H50-field-vs-coupling-transfer-lag/fig.pdf}
\caption{Field or coupling, told apart by the lag.
A simulated field moves every agent at its next call, while a coupling moves agents hop by hop, one read-out call per hop.
In the village, agents start the day within seconds of each other (a field). A peer message raises the chance of talk only from the read-out call on (a coupling).
Animation: \texttt{writeup/visuals/H50-field-vs-coupling-transfer-lag/anim.mp4}.}
\label{fig:int-H50}
\end{figure}

\textbf{What we found:}
Activity co-movement is the schedule.
The daily start and stop alone explain 0.62 (regime I) and 0.67 (regime III) of it beyond the null.
Inside the window where all agents run, the per-pair activity correlation nearly vanishes (0.056 and 0.005).
In G38 the day-start spread has an interquartile range of 9\,s, and 82\% of agents start before any peer message reaches them.
Talk co-movement is a coupling that acts only at the read-out call.
$J_1>0$ in 45/71 units and $J_1<0$ in none; the onset is at hop 1 in all 45.
Pooled, $J_1\approx0.034$ (regime I) and 0.019 (regime III), which is $+25\%$ and $+44\%$ over the base talk rate.
The talk call replaces a work call; the idle share does not change.
In regime III the coupling acts only on named messages: $J_1=0.17$ for a message that names the recipient, 0.004 otherwise.
The proposed split by regime (HH167) fails.
At NE43, switching off the nudger leaves peer coupling ($\times0.69$) and activity co-movement in place.
Scope: 71 units, 35 periods.
The claim ``activity co-movement is the scheduler's field; talk is a coupling at the hop-1 read-out call'' stands\cred{0.85}.
Its expected usefulness is EU~2.98 (value if true 3.5; two raters, one blind).
The card and \texttt{meta.json} agree.

\textbf{What it means:}
For an operator: do not read activity synchrony as influence.
First trim each day to the window in which all agents run, or stagger the starts. About 0.6 of activity synchrony is the schedule.
To spread a message in an always-on swarm, name the recipients: the jump is 0.17 with a name and 0.004 without.
Expect the response at the recipient's next call, a median of 16--24\,s, or minutes if the agent pauses.
A nudge moves the agent it targets ($J_1\approx0.14$ at hop 1) but leaves the collective mode unchanged.

\textbf{Caveats:}
Regime-I call starts are placed from latency; the jump holds on logged starts (\#24--\#31), but early regime I is unchecked.
The rising regime-I kernel beyond hop 2 has no synthetic validation.
Per-unit intervals are mildly anti-conservative (about 9\% false positives in synthetic data).
The field and coupling shares are model-based and do not add up, because coupling amplifies field-driven talk.
Room changes and roster joins confound NE43, and the \#51 relay tests have about 100 pairs.
Several numeric predictions failed: the nudge onset, the regime-I field share (0.51 against a predicted $\le0.2$) and the corner frequency.
The content channel is not done.
The fragile flag is not set.
No confirmation test on reserved goal periods has run.
